\begin{figure}[htbp]
\centering
\resizebox{0.90\textwidth}{!}{%
\begin{tikzpicture}[
  font=\sffamily\small,
  state/.style={draw=VIUDark!70,fill=VIUWhite,rounded corners=3pt,align=center,minimum height=1.1cm,text width=2.75cm,inner sep=5pt},
  event/.style={draw=VIUOrange!85!black,fill=VIUOrange!12,rounded corners=3pt,align=center,minimum height=1.1cm,text width=2.75cm,inner sep=5pt},
  flow/.style={-{Latex[length=2.5mm]},line width=0.9pt,draw=VIUDark!75}
]
\node[state] (nominal) at (0,0) {Ejecuci\'on nominal\\$C_k(t)$};
\node[event] (fault) at (3.8,0) {Evento adverso\\fallo, bater\'ia o retardo};
\node[state] (detect) at (7.6,0) {Detecci\'on y poda\\$C_k^- = C_k\setminus F$};
\node[state] (repair) at (11.4,0) {Reclutamiento local\\$C_k^+=C_k^-\cup R$};
\node[state] (guard) at (7.6,-2.25) {Guardia f\'isica\\$\rho_k\leq\varepsilon_w$};

\draw[flow] (nominal) -- (fault);
\draw[flow] (fault) -- (detect);
\draw[flow] (detect) -- (repair);
\draw[flow] (repair.south) |- (guard.east);
\draw[flow] (guard.west) -| node[pos=0.28,below,font=\scriptsize] {$t_{\rm rec}$} (nominal.south);
\end{tikzpicture}}
\caption[Ciclo de recuperaci\'on ante eventos adversos]{Ciclo de recuperaci\'on evaluado en SP6. La misi\'on solo vuelve al estado nominal cuando la coalici\'on reparada supera de nuevo la guardia de factibilidad f\'isica.}
\viuownsource
\label{fig:sp6-recovery-cycle}
\end{figure}
